% Part: incompleteness
% Chapter: representability-in-q
% Section: representable-comp

\documentclass[../../../include/open-logic-section]{subfiles}

\begin{document}

\olfileid{inc}{req}{rpc}
\olsection{在 $\Th{Q}$ 中可表示的函数都是可计算的}

我们将证明，在~$\Th{Q}$ 中可表示的函数都是可计算的。为此，首先需要证明一个关于在~$\Th{Q}$ 中可表示的函数的引理。

\begin{lem}\ollabel{lem:rep-q}
  如果 $f(x_0,
  \dots, x_k)$ 在~$\Th{Q}$ 中可表示，则存在公式~$!A(x_0, \dots, x_k, y)$，使得
  \[
  \Th{Q} \Proves !A_f(\num{n_0}, \dots, \num{n_k}, \num{m})
  \quad \text{当且仅当} \quad m = f(n_0, \dots, n_k).
  \]
\end{lem}

\begin{proof}
  “当”的方向由
\olref[int]{defn:representable-fn}\olref[int]{defn:rep:a} 给出。“仅当”的方向论证如下：假设 $\Th{Q} \Proves !A_f(\num{n_0},
\dots, \num{n_k}, \num{m})$ 但 $m \neq f(n_0, \dots, n_k)$。令 $l =
f(n_0, \dots, n_k)$。根据
\olref[int]{defn:representable-fn}\olref[int]{defn:rep:a}，$\Th{Q}
\Proves !A_f(\num{n_0}, \dots, \num{n_k}, \num{l})$。根据
\olref[int]{defn:representable-fn}\olref[int]{defn:rep:b}，
$\lforall[y][(!A_f(\num{n_0}, \dots, \num{n_k}, y) \lif \num{l} =
y)]$。根据逻辑及 $\Th{Q} \Proves
!A_f(\num{n_0}, \dots, \num{n_k}, \num{m})$ 的假设，可得 $\Th{Q}
\Proves \eq[\num{l}][\num{m}]$。另一方面，由
\olref[bre]{lem:q-proves-neq}，$\Th{Q} \Proves
\eq/[\num{l}][\num{m}]$。于是 $\Th{Q}$ 不一致。但这是不可能的，因为 $\Th{Q}$ 被标准模型所满足（参见
\olref[int][def]{def:standard-model}），即 $\Sat{N}{\Th{Q}}$，而根据可靠性定理，可满足的理论总是一致的（\tagrefs{prfAX/{fol:axd:sou:cor:consistency-soundness},
prfSC/{fol:seq:sou:cor:consistency-soundness},
prfND/{fol:ntd:sou:cor:consistency-soundness},
prfTab/{fol:tab:sou:cor:consistency-soundness}}）。
\end{proof}

\begin{lem}
每个在 $\Th{Q}$ 中可表示的函数都是可计算的。
\end{lem}

\begin{proof}
我们先直观地说明为什么这是成立的。为了计算~$f$，我们如下进行：枚举算术语言中所有可能的推导~$\delta$。这是可以机械完成的。对于每个推导，检查它是否是形如~$!A_f(\num{n_0}, \dots,
\num{n_k}, \num{m})$（即根据 \olref{lem:rep-q} 在~$\Th{Q}$ 中表示 $f$ 的公式）的公式的一个推导。如果是，则由 \olref{lem:rep-q}，$m = f(n_0, \dots, n_k)$，于是我们找到了~$f$ 的值。由于 $\Th{Q} \Proves !A_f(\num{n_0}, \dots, \num{n_k},
\num{f(n_0, \dots, n_k)})$，搜索必定终止，即我们最终会找到一个符合要求的 $\delta$。

这并不完全精确，因为我们的过程处理的对象是推导和公式，而不仅仅是数字，并且我们还没有确切解释为什么“枚举所有可能的推导”是机械可行的。但正如我们所见，可以用哥德尔数来编码项、公式和推导。我们也引入了一个精确的计算模型，即一般递归函数。而且我们已经看到，关系 $\Prf[\Th{Q}](d,y)$ 是（原始）递归的，该关系成立当且仅当 $d$ 是这样一个推导的哥德尔数：该推导是从~$\Th{Q}$ 的公理出发对哥德尔数为~$y$ 的公式所作的推导。我们需要的其他原始递归函数还有 $\fn{num}$（\olref[art][trm]{prop:num-primrec}）和 $\fn{Subst}$（\olref[art][sub]{prop:subst-primrec}）。利用这些函数，可以通过极小化来定义~$f$；因此 $f$ 是递归的。

首先，定义
\begin{multline*}
  A(n_0, \dots, n_k, m) = \\
  \fn{Subst}(\fn{Subst}(\dots\fn{Subst}(\Gn{!A_f}, \fn{num}(n_0), \Gn{x_0}),\\ \dots),
  \fn{num}(n_k),  \Gn{x_k}), \fn{num}(m), \Gn{y})
\end{multline*}
这看起来有些复杂，但它其实就是函数 $A(n_0, \dots, n_k,
m) = \Gn{!A_f(\num{n_0}, \dots, \num{n_k}, \num{m})}$。

现在，考虑关系~$R(n_0, \dots, n_k, s)$，它成立当且仅当 $(s)_0$ 是从~$\Th{Q}$ 出发对 $!A_f(\num{n_0}, \dots, \num{n_k}, \num{(s)_1})$ 的一个推导的哥德尔数：
\[
R(n_0, \dots, n_k, s) \quad \text{当且仅当} \quad \Prf[\Th{Q}]((s)_0, A(n_0,
\dots, n_k, (s)_1))
\]
如果我们能找到一个使 $R(n_0, \dots, n_k, s)$ 成立的~$s$，我们就找到了一对数——$(s)_0$ 和~$(s)_1$——使得 $(s)_0$ 是 $A_f(\num{n_0}, \dots,
\num{n_k}, (s)_1)$ 的一个推导的哥德尔数。因此，寻找这样的~$s$ 就好比在非形式化证明中寻找数对 $d$ 和 $m$。而一个“寻找”这样的 $s$ 的可计算函数可以通过正则极小化来定义。注意 $R$ 是正则的：对每个 $n_0$, \dots, $n_k$，都存在 $\Th{Q} \Proves !A_f(\num{n_0}, \dots,
\num{n_k}, \num{f(n_0, \dots, n_k)})$ 的一个推导~$\delta$，因此对于 $s = \tuple{\Gn{\delta}, f(n_0, \dots, n_k)}$，$R(n_0, \dots, n_k, s)$ 成立。所以，我们可以将 $f$ 写作
\[
f(n_0,\dots,n_{k}) = (\umin{s}{R(n_0, \dots, n_k, s)})_1.
\]
\end{proof}

\end{document}